\documentclass[a4paper]{article}

\usepackage[english]{babel}
%\usepackage[utf8x]{inputenc}
\usepackage{amsmath}
\usepackage{graphicx}
\usepackage[colorinlistoftodos]{todonotes}
\usepackage[all]{xy}
\usepackage{dsfont}
\usepackage{amssymb}
\usepackage{bbold}
\usepackage{amsthm}

\usepackage{tikz}
\usepackage{verbatim}
\usetikzlibrary{matrix}
\usetikzlibrary{cd}

\usepackage[autostyle=true]{csquotes}
\usepackage{pdfpages}
\usepackage{hyperref}
\usepackage{fancyhdr}



\newtheorem{theorem}{Theorem}[section]
\newtheorem{proposition}{Proposition}[theorem]
\newtheorem{lemma}{Lemma}[theorem]
\newtheorem{corollary}{Corollary}[theorem]
\newtheorem*{definition}{Definition}


\newcommand{\norm}[1]{\left\lVert#1\right\rVert}

\newcommand{\Rn}{\mathbb{R}^n}


\newcommand{\too}[1]{\overset{#1}{\longrightarrow}}
\newcommand{\hookar}[1]{\overset{#1}{\hookrightarrow}}
\newcommand{\colim}[1]{ \underset{#1}{\texttt{colim}} \text{ } }
\newcommand{\fset}[1]{ \{ #1_0 , \dots , #1_n \} }
\newcommand{\ddt}{\frac{d}{dt}}


\newcommand{\N}{\mathbb{N}}
\newcommand{\Z}{\mathbb{Z}}







\fancyhead[L]{\includegraphics[scale=0.1]{UFT.png}}

\begin{document}


\thispagestyle{fancy}
\vspace*{1cm}

\begin{minipage}[r]{12cm}
    \begin{flushright}
\begin{tikzpicture}
\draw[green, very thick] (0,0) rectangle (3,2);
\end{tikzpicture}
    \end{flushright}
\end{minipage}

\begin{center}
\Large{\textbf{Control 2\\ Pauta\\ Álgebra Lineal\\ }}
\end{center}
\vspace{1cm}
\begin{minipage}[r]{15cm}
\begin{flushleft}



\end{flushleft}
\end{minipage}
\vspace{2cm}
\begin{center}
\Large{\textbf{Facultad de Ingeniería}}
\end{center}
\vspace{1cm}
\textbf{En caso de que un alumno sea sorprendido en actitudes deshonestas, ya sea intentando o
dejándose copiar, será calificado con nota mínima 1.0 y perderá toda posibilidad de eliminar
dicha nota. Además, esta conducta quedará archivada como antecedente en su carpeta y
podrá ser causal de eliminación.
Así mismo, está prohibido salir de la sala, usar o mantener sobre la mesa teléfonos celulares
o cualquier tipo de aparato electrónico, a no ser que el profesor lo requiera.}\\



\vspace{5cm}
\makebox[12cm]{\textbf{Nombre Alumno}:\ \hrulefill} 
\newpage
\thispagestyle{empty}
\section*{Preguntas}
\begin{enumerate}


    \item \textbf{(12 puntos)} Dado el siguiente sistema de ecuaciones lineales:

\[
\begin{aligned}
2x + 4y + (c^2 - 1)z &= c + 1 \\
x + y + z &= 1 \\
x + 3y + 2z &= 2
\end{aligned}
\]

Determine los valores de \( c \in \mathbb{R} \) de modo que el sistema:

\begin{itemize}
    \item Tenga solución única
    \item Tenga infinitas soluciones
    \item No tenga solución
\end{itemize}

\textbf{Solución: }

Escribimos la matriz aumentada del sistema:

\[
\left[
\begin{array}{ccc|c}
2 & 4 & c^2 - 1 & c + 1 \\
1 & 1 & 1 & 1 \\
1 & 3 & 2 & 2
\end{array}
\right]
\]

Aplicamos operaciones elementales:

\begin{itemize}
    \item \( f_1 - 2f_2 \rightarrow f_1 \)
    \item \( f_3 - f_2\rightarrow f_3  \)
\end{itemize}

\[
\left[
\begin{array}{ccc|c}
0 & 2 & c^2 - 3 & c - 1 \\
1 & 1 & 1 & 1 \\
0 & 2 & 1 & 1
\end{array}
\right]
\]

Ahora restamos la nueva fila 1 a la fila 3:

\[
\left[
\begin{array}{ccc|c}
0 & 2 & c^2 - 3 & c - 1 \\
1 & 1 & 1 & 1 \\
0 & 0 & 4 - c^2 & 2 - c
\end{array}
\right]=\left[
\begin{array}{ccc|c}
0 & 2 & c^2 - 3 & c - 1 \\
1 & 1 & 1 & 1 \\
0 & 0 & (2-c)(2+c) & 2 - c
\end{array}
\right]
\]

A partir de aquí se analizan los valores de \( c \) según el sistema resultante. Como se obtiene la fila:

\[
0x + 0y + (4 - c^2)z = 2 - c
\]

Analizamos los siguientes casos:

\begin{itemize}
    \item \textbf{Solución única:} ocurre cuando \( 4 - c^2 \ne 0 \Rightarrow c \ne \pm 2 \).
    \item \textbf{Infinitas soluciones:} ocurre cuando \( 4 - c^2 = 0 \) y \( 2 - c = 0 \Rightarrow c = 2 \).
    \item \textbf{No tiene solución:} ocurre cuando \( 4 - c^2 = 0 \) y \( 2 - c \ne 0 \Rightarrow c = -2 \).
\end{itemize}
\newpage
\item
\begin{enumerate}
    \item[(a)] \textbf{(6 puntos)} Determine los valores de \( k \in \mathbb{R} \) que hagan que el conjunto
    \[
\left\{
\begin{pmatrix} 1 \\ k \\ 0 \end{pmatrix},\ 
\begin{pmatrix} -1 \\ 2 \\ 1 \end{pmatrix},\ 
\begin{pmatrix} 2 \\ -1 \\ k \end{pmatrix}
\right\}
\]

sea linealmente independiente.

\vspace{1em}

\textbf{Solución:}

Un conjunto de tres vectores en \( \mathbb{R}^3 \) es linealmente independiente si el determinante de la matriz que forman como columnas (o filas) es distinto de cero.

Formamos la matriz con estos vectores como columnas:

\[
A = \begin{bmatrix}
1 & -1 & 2 \\
k & 2 & -1 \\
0 & 1 & k
\end{bmatrix}
\]

Calculamos el determinante:

\[
\det(A) =
\begin{vmatrix}
1 & -1 & 2 \\
k & 2 & -1 \\
0 & 1 & k
\end{vmatrix}=k^2+4k+1
\]




El conjunto es linealmente independiente si:

\[
\det(A) \ne 0 \Rightarrow k^2+4k+1 \ne 0 \Rightarrow k \ne -2\pm \sqrt{3}
\]

\vspace{1em}

\textbf{Conclusión:}  
El conjunto es linealmente independiente si y solo si:
\[
\boxed{k \ne -2\pm \sqrt{3}}
\]
    
    \vspace{1em}
    
    \item[(b)] \textbf{(6 puntos)} Escriba el vector 
    \[
    \vec{v} = 3x^2 - 9x + 11 \in \mathcal{P}_2(\mathbb{R})
    \]
    como combinación lineal de los vectores 
    \[
    \vec{v}_1 = -x + 1,\quad \vec{v}_2 = x^2 + 2x,\quad \vec{v}_3 = 3x^2 - x + 5.
    \]
    
    \textbf{Solución:}
    
    Queremos encontrar escalares \(a, b, c \in \mathbb{R}\) tales que:
    \[
    \vec{v} = a\vec{v}_1 + b\vec{v}_2 + c\vec{v}_3
    \]
    
    Sustituimos:
    \[
    3x^2 - 9x + 11 = a(-x + 1) + b(x^2 + 2x) + c(3x^2 - x + 5)
    \]
    
    Expandimos:
    \[
    a(-x + 1) = -a x + a \\
    b(x^2 + 2x) = b x^2 + 2b x \\
    c(3x^2 - x + 5) = 3c x^2 - c x + 5c
    \]
    
    Sumamos todo:
    \[
    (b + 3c)x^2 + (-a + 2b - c)x + (a + 5c)
    \]
    
    Igualamos coeficientes con \(3x^2 - 9x + 11\):

    \[
    \begin{aligned}
    b + 3c &= 3 \\
    -a + 2b - c &= -9 \\
    a + 5c &= 11
    \end{aligned}
    \]

    Resolviendo el sistema:

    1. De la primera ecuación: \(b = 3 - 3c\)

    2. Sustituimos en la segunda:
    \[
    -a + 2(3 - 3c) - c = -9 \Rightarrow -a + 6 - 6c - c = -9 \Rightarrow -a - 7c = -15 \Rightarrow a = 15 - 7c
    \]

    3. Sustituimos \(a\) en la tercera:
    \[
    15 - 7c + 5c = 11 \Rightarrow -2c = -4 \Rightarrow c = 2
    \]

    Luego:
    \[
    b = 3 - 3(2) = -3,\quad a = 15 - 7(2) = 1
    \]

    \textbf{Conclusión:}
    \[
    \vec{v} = a \vec{v}_1 + b \vec{v}_2 + c \vec{v}_3 = \vec{v}_1 - 3\vec{v}_2 + 2\vec{v}_3
    \]
\end{enumerate}

\newpage
\item

Encuentre en cada caso una base y la dimensión del siguiente espacio vectorial:

\[
V = \left\langle
\begin{pmatrix} 0 & -1 \\ 1 & 1 \end{pmatrix},
\begin{pmatrix} 0 & 1 \\ -2 & 3 \end{pmatrix},
\begin{pmatrix} 0 & 3 \\ 1 & 4 \end{pmatrix},
\begin{pmatrix} 0 & 2 \\ -4 & 6 \end{pmatrix}
\right\rangle
\]

\textbf{Solución:}

\[
V = \left\langle
\begin{pmatrix} 0 & -1 \\ 1 & 1 \end{pmatrix},\ 
\begin{pmatrix} 0 & 1 \\ -2 & 3 \end{pmatrix},\ 
\begin{pmatrix} 0 & 3 \\ 1 & 4 \end{pmatrix},\ 
\begin{pmatrix} 0 & 2 \\ -4 & 6 \end{pmatrix}
\right\rangle
\]

\textbf{Solución:}

Representamos cada matriz como un vector en \(\mathbb{R}^4\), leyendo por filas:

\[
A_1 = \begin{bmatrix} 0 \\ -1 \\ 1 \\ 1 \end{bmatrix}, \quad
A_2 = \begin{bmatrix} 0 \\ 1 \\ -2 \\ 3 \end{bmatrix}, \quad
A_3 = \begin{bmatrix} 0 \\ 3 \\ 1 \\ 4 \end{bmatrix}, \quad
A_4 = \begin{bmatrix} 0 \\ 2 \\ -4 \\ 6 \end{bmatrix}
\]

Si el determinante es distinto de cero, los vectores son linealmente independientes y forman una base del subespacio.

Consideramos la siguiente matriz:

\[
M = \begin{bmatrix}
0 & 0 & 0 & 0 \\
-1 & 1 & 3 & 2 \\
1 & -2 & 1 & -4 \\
1 & 3 & 4 & 6
\end{bmatrix}
\]

Cuyo determinante es claramente cero. Eliminamos el \'ultimo vector y haciendo reducci\'on por filas obtenemos:

\[
M' = \begin{bmatrix}
0 & 0 & 0 \\
-1 & 1 & 3  \\
1 & -2 & 1  \\
1 & 3 & 4 
\end{bmatrix}\simeq \begin{bmatrix}
 1 & 0 & 0 \\
 0& 1 & 0 \\
 0 & 0 & 1 \\
0 & 0 & 0
\end{bmatrix}
\]


\textbf{Conclusión:}

\begin{itemize}
\item Una base del espacio vectorial es:
\[
\left\{
\begin{pmatrix} 0 & -1 \\ 1 & 1 \end{pmatrix},\ 
\begin{pmatrix} 0 & 1 \\ -2 & 3 \end{pmatrix},\ 
\begin{pmatrix} 0 & 3 \\ 1 & 4 \end{pmatrix}
\right\}
\]
\item La dimensión del espacio vectorial \(V\) es 3.
\end{itemize}
La matriz $\begin{pmatrix} 1 & 0 \\0 & 0 \end{pmatrix}$ no pertenece a $V$, pues haciendo la combinación lineal, se obtiene que 
\[
\begin{pmatrix} 1 & 0 \\0 & 0 \end{pmatrix}=\alpha
\begin{pmatrix} 0 & -1 \\ 1 & 1 \end{pmatrix}+\beta\ 
\begin{pmatrix} 0 & 1 \\ -2 & 3 \end{pmatrix}+\gamma\ 
\begin{pmatrix} 0 & 3 \\ 1 & 4 \end{pmatrix}
\]
resulta $1=0\alpha+0\beta+0\gamma=0$, lo cual es imposible.
\newpage


\item Muestre que el conjunto

\[
W = \left\{
\begin{pmatrix}
x \\
y \\
z
\end{pmatrix}
\in \mathbb{R}^3 \;\middle|\; 2x - y = 0
\right\}
\]

es subespacio vectorial de \(\mathbb{R}^3\). Determine la dimensión de este espacio.

\subsection*{Solución}

Para que \(W\) sea subespacio vectorial de \(\mathbb{R}^3\), debe cumplir las siguientes condiciones:

\begin{itemize}
  \item Contener al vector cero.
  \item Ser cerrado bajo suma de vectores.
  \item Ser cerrado bajo multiplicación por escalar.
\end{itemize}



\[
\begin{pmatrix} 0 \\ 0 \\ 0 \end{pmatrix} \in \mathbb{R}^3,\quad \text{y se verifica que } 2(0) - 0 = 0 \Rightarrow \text{pertenece a } W.
\]


Sean \(\vec{u} = \begin{pmatrix} x_1 \\ y_1 \\ z_1 \end{pmatrix}, \vec{v} = \begin{pmatrix} x_2 \\ y_2 \\ z_2 \end{pmatrix} \in W\), es decir, \(2x_1 - y_1 = 0\) y \(2x_2 - y_2 = 0\). Entonces:

\[
\vec{u} + \vec{v} = \begin{pmatrix} x_1 + x_2 \\ y_1 + y_2 \\ z_1 + z_2 \end{pmatrix},\quad
2(x_1 + x_2) - (y_1 + y_2) = (2x_1 - y_1) + (2x_2 - y_2) = 0 + 0 = 0
\]

Entonces \(\vec{u} + \vec{v} \in W\).



Sea \(\vec{u} = \begin{pmatrix} x \\ y \\ z \end{pmatrix} \in W\) tal que \(2x - y = 0\), y \(\lambda \in \mathbb{R}\). Entonces:

\[
\lambda \vec{u} = \begin{pmatrix} \lambda x \\ \lambda y \\ \lambda z \end{pmatrix},\quad
2(\lambda x) - (\lambda y) = \lambda (2x - y) = \lambda \cdot 0 = 0
\]

Por tanto, \(\lambda \vec{u} \in W\).

\bigskip

\textbf{Conclusión:} \(W\) es subespacio vectorial de \(\mathbb{R}^3\).

Cálculo de la dimensión: La condición \(2x - y = 0\) equivale a:

\[
y = 2x
\]

Entonces, todo vector de \(W\) se puede escribir como:

\[
\begin{pmatrix}
x \\
2x \\
z
\end{pmatrix}
= x \begin{pmatrix} 1 \\ 2 \\ 0 \end{pmatrix}
+ z \begin{pmatrix} 0 \\ 0 \\ 1 \end{pmatrix}
\]

Por tanto, una base de \(W\) es:

\[
\left\{
\begin{pmatrix} 1 \\ 2 \\ 0 \end{pmatrix},\ 
\begin{pmatrix} 0 \\ 0 \\ 1 \end{pmatrix}
\right\}
\]

\textbf{Y la dimensión de \(W\) es: 2.}

\newpage

\item Considere la matriz
\[
A = 
\begin{pmatrix}
0 & 3 & 0 \\
1 & -2 & 0 \\
2 & 1 & -1
\end{pmatrix}
\]

\begin{itemize}
\item[(a)] Obtenga los valores propios de \(A\) y sus vectores propios asociados.
\item[(b)] Presente, si es posible, la diagonalización de esta matriz, encontrando matrices \(P\) y \(D\) que cumplan
\[
A = P D P^{-1}
\]
\end{itemize}

\subsection*{Solución}

\textbf{Paso 1: Calcular los valores propios}

Buscamos los valores propios resolviendo el polinomio característico:

\[
\det(A - \lambda I) = 0
\]
Expandimos por la tercera columna:
\begin{align*}
    \det
\begin{pmatrix}
-\lambda & 3 & 0 \\
1 & -2 - \lambda & 0 \\
2 & 1 & -1 - \lambda
\end{pmatrix}
&= 0\\
%&= (-\lambda)
% \left|
% \begin{matrix}
% -2 - \lambda & 0 \\
% 1 & -1 - \lambda
% \end{matrix}
% \right|\\
&= (-1-\lambda)[(- \lambda)(-2 - \lambda)-3]\\
&= (-1-\lambda)[\lambda^2 + 2\lambda-3]\\
&= (-1-\lambda)[(\lambda+3)(\lambda-1)]
\end{align*}

Los valores propios son:
\[
\boxed{\lambda_1 = -3, \quad \lambda_2 = -1, \quad \lambda_3 = 1}
\]

Los vectores propios asociados son:

\[
\lambda = -3 \quad \Rightarrow \quad 
\vec{v}_1 = 
\begin{pmatrix}
-2 \\
2 \\
1
\end{pmatrix}
\]

\[
\lambda = -1 \quad \Rightarrow \quad 
\vec{v}_2 = 
\begin{pmatrix}
0 \\
0 \\
1
\end{pmatrix}
\]

\[
\lambda = 1 \quad \Rightarrow \quad 
\vec{v}_3 = 
\begin{pmatrix}
6 \\
2 \\
7
\end{pmatrix}
\]

\textbf{(b) Diagonalización}

Dado que \(A\) tiene 3 valores propios distintos con 3 vectores propios linealmente independientes, se puede diagonalizar.

La matriz \(P\) se forma con los vectores propios como columnas, y la matriz \(D\) es diagonal con los valores propios:

\[
P = \begin{pmatrix}
-2 & 0 & 6 \\
2 & 0 & 2 \\
1 & 1 & 7
\end{pmatrix}, \quad
D = \begin{pmatrix}
-3 & 0 & 0 \\
0 & -1 & 0 \\
0 & 0 & 1
\end{pmatrix}
\]

Entonces, la diagonalización es:
\[
\boxed{A = P D P^{-1}}
\]
\end{enumerate}


\end{document}